\subsection*{Study A: with Claude Opus 5.5 every executor reached the correct verdict with or without the specification, and generic reviews turned three correct verdicts wrong}
All 40 executor runs and all 120 final deliverables were completed and
graded. Every executor run reached the key's verdict, in both arms (20 of
20 with the specification and 20 of 20 without it; risk difference 0,
Newcombe 95\% interval $-16$ to $+16$ percentage points; Fisher
$p = 1$; Mantel--Haenszel $p = 1$) (Fig~\ref{fig:studyA}A). After review
and repair, 40 of 40 deliverables were correct in the checklist arm and 37
of 40 in the generic-review arm (exact McNemar $p = 0.25$ against the
checklist; $p = 0.25$ against no review). None of the four primary
contrasts was significant (Holm-adjusted $p = 1$ for each). The two
graders agreed on the verdict of all 120 deliverables.

\begin{figure}[H]
\caption{\textbf{Study A: controlled execution study with Claude Opus~5.5
as executor, reviewer and repairer.} (A)~Number of the five executor runs
per task whose final deliverable reached the key's verdict, by executor
arm (paper: brief, data and source paper; contract: the same plus the
deployed specification) and by what followed the executor (no review, a
generic critical review then repair, or a generic-checklist review then
repair). (B)~Share of each task's known pitfalls handled correctly (mean
over deliverables; line, 95\% bootstrap interval over executor runs
pooled over tasks).
(C)~Review points judged valid or invalid (or unclear) by a grader with
the key, over the 40 reviews of each style.}
\label{fig:studyA}
\end{figure}

The three verdicts that became wrong were all on the cross-model
alignment task (T4) in the contract arm (where the executor also received
the deployed specification), after a generic review. In each,
the reviewer argued that the executor's decision rule was too lenient, or
that metrics named in the specification (the mean in-sample canonical
correlation, or the Spearman correlation of pairwise distances) gave a
weaker result than the executor's chosen metric. The repairing agent
accepted these points, rewrote its decision rule, and changed the verdict
from ``aligned with scGPT but clearly less than with Geneformer'' to
``inconclusive''. The key does not accept ``inconclusive'': every estimate
is far above chance, and the Geneformer-minus-scGPT difference has an
interval that excludes zero on every panel. Canonical correlation finds
pairs of directions, one in each embedding table, along which the genes'
positions correlate as strongly as possible. When it is fitted and scored
on the same genes (in-sample), as the reviewers used it, it is high even
for shuffled genes.
In the setting these executors used (20 principal components, mean of all
20 canonical correlations), their own permutation tests put chance at
about 0.20. The observed values were 0.40 for scGPT and 0.44 for
Geneformer. This chance level depends on the number of components as well
as on the number of genes, and the specification does not state it.
Graders with the key judged 17 of the 22 points in these three reviews
invalid. The checklist reviews of the same three executor outputs
did not change the verdict.

The secondary outcomes moved little (Fig~\ref{fig:studyA}B). Executors
handled most of each task's pitfalls without help (94.5\% on average
without review). Review and repair raised this slightly: by 2.8
percentage points with the checklist (95\% bootstrap interval 0.9 to 4.7)
and 2.0 points with a generic review (0.3 to 3.7); the checklist and the
generic review did not differ ($+0.8$; $-0.9$ to $+2.5$). Without review,
executors with the specification handled 2.3 points fewer pitfalls than
those without it ($-5.4$ to $+0.8$) and reported 9.2 points more key
numbers within tolerance ($-0.6$ to $+21.3$). These intervals resample
executor runs pooled over tasks. The pre-registered bootstrap resamples
runs within each task. It gives the same conclusions for the effects of
review on pitfalls handled, in percentage points: 1.1 to 4.4 for the
checklist against no review, 0.5 to 3.4 for a generic review against no
review, and $-0.9$ to $+2.4$ for the checklist against a generic review.
It also gives the same conclusions for the effects of review on pitfalls
handled in the Haiku arm below. But for the two differences between the
executor arms, its intervals just exclude zero: $-4.4$ to $-0.2$ points
for pitfalls handled and $+0.6$ to $+18.9$ points for key numbers within
tolerance. Both
differences are small, and whether their intervals exclude zero depends on
how the runs are resampled. Materially false statements were rare (0.13
per deliverable) and did not differ between conditions.

Most review points were judged invalid by a grader with the key, meaning
the point was wrong, already handled in the deliverable, or did not
matter for its results: 61\% of the 281 generic-review points and 70\% of
the 297 checklist points (208 invalid and 1 unclear)
(Fig~\ref{fig:studyA}C). Only in the three T4 cases above did acting on
such points turn a correct verdict into a wrong one. Three other repairs,
all on T1 in the paper arm (two after a generic review and one after a
checklist review), changed ``not supported'' to ``inconclusive''. The T1
key accepts that answer under stated conditions, so all three stayed
correct.

An executor run took a median of 27~minutes and 27 tool calls; a review
took a median of about 14~minutes and 16 (checklist) or 18.5 (generic)
tool calls, and a repair 21 to 25 tool calls.
All 120 returned analysis scripts ran on the task data and printed their
results. Six of them were stopped at least once because the shared
machine was busy: five ran past the 20-minute limit, and two were stopped
without output (one script had both). Run again with fewer scripts at a
time, or on their own, all six finished in under 10~minutes. Of the
11,045 estimates that the deliverables reported, 10,363 (94\%) could be matched by name to a number
the script printed, and all 10,363 agreed within tolerance. Of the 682
unmatched estimates, 681 were in T3, whose scripts printed their results
under other names. For these, the only check left was to look for each
value among all the numbers a script printed, within the same tolerance,
and that check says little. In the 21 T3 scripts concerned, 98--100\% of
the deliverable's own estimates were found this way, but so were
61--97\% (median 89\%) of the same number of estimates drawn at random
from deliverables of other tasks (mean of 50 draws per script).

In the exploratory arm with Claude Haiku~4.5 as executor, reviewer and
repairer (S2~Appendix), executors without review reached the key's verdict
in 15 of 20 runs with the specification and in 10 of 20 without it
(difference $+25$ points, Newcombe 95\% interval $-5$ to $+49$; Fisher
$p = 0.19$); most of the difference came from T1 (3 of 5 against 0 of 5).
Haiku executors handled far fewer pitfalls than Opus executors (30--36\%
against 93--96\% without review). A checklist review followed by repair
raised the share of pitfalls handled by 9.2 percentage points (95\%
bootstrap interval over executor runs pooled over tasks, 4.9 to 13.6;
pre-registered interval resampling runs within each task, 5.0 to 13.3).
It raised it by 6.4 points more than a generic review followed by repair
(over the 28 runs with both reviews valid: 1.3 to 11.7 pooled over tasks,
2.1 to 10.4 within tasks). But it did not
raise the share of correct verdicts. Of the 33 executor runs with a valid
checklist review, repairs turned 3 of 12 wrong verdicts right and 3 of 21
right verdicts wrong. Of the 35 executor runs with a valid generic review,
repairs turned 2 of 13 wrong verdicts right and none of 22 right verdicts
wrong. Twelve reviews (7 checklist, 5 generic) returned no valid output,
so 12 of the 120 deliverables are missing.

